% =============================================================================
% Smart Grid MCP — Data Pipeline Documentation
% =============================================================================
% Describes how the five Kaggle datasets are downloaded, reconciled on a
% synthetic transformer_id key, and processed into the six CSVs consumed
% by the four MCP servers.
%
% Intended audience: team members, paper reviewers, open-source contributors.
% Suitable for inclusion as an appendix in the NeurIPS / IEEE final report.
% =============================================================================

\section{Data Pipeline}
\label{sec:data-pipeline}

\subsection{Source Datasets}
\label{sec:datasets}

We draw from five publicly available Kaggle datasets covering the four
AssetOpsBench agent domains (IoT, FMSR, TSFM, WO) for power transformers.
Table~\ref{tab:datasets} summarises their properties.

\begin{table}[h]
\centering
\caption{Source datasets used in the Smart Grid benchmark extension.}
\label{tab:datasets}
\begin{tabular}{lllll}
\hline
\textbf{Dataset} & \textbf{Rows} & \textbf{Domain} & \textbf{License} & \textbf{Key challenge} \\
\hline
Power Transformers FDD \& RUL~\cite{fddrul}
  & 3{,}000 files $\times$ 420 & TSFM, IoT & CC0 & Normalised; no timestamps \\
DGA Fault Classification~\cite{dga}
  & 201 & FMSR & CC0 & No asset ID \\
Transformer Health Index~\cite{healthindex}
  & 470 & FMSR & ODbL & No asset ID \\
Current \& Voltage Monitoring~\cite{monitoring}
  & 19{,}352 & IoT, TSFM & \copyright\ Authors & Single transformer \\
Smart Grid Fault Records~\cite{faultrecords}
  & 506 & WO & CC0 & No asset ID \\
\hline
\end{tabular}
\end{table}

\noindent
Only datasets licensed CC0 will be included in the open-source pull request
to AssetOpsBench.  The ODbL (Open Database License) and author-copyright
datasets are used locally for development and will not be redistributed.

\subsubsection{FDD \& RUL Dataset}

The Power Transformers FDD \& RUL dataset contains 3{,}000 per-transformer
time-series files, each with 420 time steps and four dissolved-gas columns
(\texttt{H2}, \texttt{CO}, \texttt{C2H4}, \texttt{C2H2}).  Values are
normalised gas-concentration ratios rather than raw ppm measurements.

Separate label files (\texttt{labels\_fdd\_train.csv},
\texttt{labels\_rul\_train.csv}) provide per-file fault category and remaining
useful life (RUL) estimates.  The fault category is an integer in $\{1, 2, 3,
4\}$, where 1 denotes normal operation and 4 denotes a serious fault.  RUL
values range from 362 to 1{,}093 (units: days).  The training split contains
2{,}100 files; the test split contains 900.

\subsubsection{DGA Fault Classification Dataset}

Contains 201 dissolved-gas analysis (DGA) records with raw gas concentrations
in ppm (\texttt{H2}, \texttt{CH4}, \texttt{C2H6}, \texttt{C2H4},
\texttt{C2H2}) and a textual fault-type label.  The seven observed fault types
are: \emph{Arc discharge}, \emph{Spark discharge},
\emph{High-temperature overheating}, \emph{Low-temperature overheating},
\emph{Partial discharge}, \emph{Low/Middle-temperature overheating}, and
\emph{Middle-temperature overheating}.  All 201 records represent faulty
transformers; there are no normal-operation examples.

\subsubsection{Transformer Health Index Dataset}

Contains 470 records with expanded gas columns (additionally \texttt{Nitrogen},
\texttt{Oxigen}, \texttt{DBDS}, \texttt{Power factor}, \texttt{Interfacial V},
\texttt{Dielectric rigidity}, \texttt{Water content}) and a scalar
\texttt{Health index} (range 13.4–95.2; higher is healthier).  This dataset is
used to source DGA records for healthy transformers, since the DGA Fault
Classification dataset contains only faulty examples.

\subsubsection{Current \& Voltage Monitoring Dataset}

Five CSV files capturing real time-series measurements from a single
transformer over approximately 2019–2020.  Relevant tables:
\begin{itemize}
  \item \texttt{Overview.csv} (20{,}316 rows): oil temperature index (OTI),
        winding temperature index (WTI), ambient temperature (ATI), oil level
        (OLI).
  \item \texttt{CurrentVoltage.csv} (19{,}352 rows): per-phase voltages
        (VL1, VL2, VL3) and currents (IL1, IL2, IL3) in V and A.
  \item \texttt{Power.csv}, \texttt{PowerFactor.csv}, \texttt{TotalPower.csv}:
        power, THD, and energy metrics.
\end{itemize}

\subsubsection{Smart Grid Fault Records Dataset}

Contains 506 fault event records with attributes including fault type, location
(latitude/longitude), voltage, current, power load, temperature, weather
condition, maintenance status, component health, fault duration, and downtime.
This dataset has no asset identifier and covers grid-level events rather than
individual transformer units.

\subsection{The Shared-Key Problem}
\label{sec:shared-key}

A fundamental challenge in assembling this benchmark is that the five source
datasets are entirely \emph{disjoint}: they originate from different research
groups, cover different physical assets, and share no common identifier.  The
AssetOpsBench multi-domain scenario format requires that a single agent call
tools from multiple domains (IoT, FMSR, TSFM, WO) for the \emph{same} asset
within one task trajectory.  This is impossible without a join key.

We resolve this by synthesising a \texttt{transformer\_id} key (\texttt{T-001}
through \texttt{T-020}) and assigning rows from each source dataset to
transformer IDs using the principled strategy described in
Section~\ref{sec:id-strategy}.  This approach is consistent with how
AssetOpsBench handles other domains where data from multiple disparate sources
must be reconciled \cite{assetopsbench}.

\subsection{Transformer Selection Strategy}
\label{sec:id-strategy}

We select 20 representative transformers from the FDD \& RUL training set,
stratified across four health tiers, to ensure the benchmark covers the full
range of operating conditions an agent may encounter in production:

\begin{table}[h]
\centering
\caption{Transformer health tiers and selection criteria.}
\label{tab:tiers}
\begin{tabular}{lllll}
\hline
\textbf{Tier} & \textbf{IDs} & \textbf{FDD Cat.} & \textbf{RUL} & \textbf{Count} \\
\hline
Healthy (long life)  & T-001 – T-005 & 1 & $\geq$ 1{,}093 days & 5 \\
Healthy (aging)      & T-006 – T-010 & 1 & $\leq$ 560 days & 5 \\
Minor/moderate fault & T-011 – T-015 & 2–3 & mixed & 5 \\
Serious fault        & T-016 – T-020 & 4 & mixed & 5 \\
\hline
\end{tabular}
\end{table}

\noindent
Each transformer ID is anchored to one FDD \& RUL training file, which provides
420 time steps of dissolved-gas time-series data and an authoritative RUL label.

\subsection{Processed Output Files}
\label{sec:outputs}

The pipeline (\texttt{data/build\_processed.py}) produces six CSV files in
\texttt{data/processed/}, one per MCP server domain.  Table~\ref{tab:outputs}
summarises them; detailed schemas follow.

\begin{table}[h]
\centering
\caption{Processed output files and their MCP server consumers.}
\label{tab:outputs}
\begin{tabular}{llll}
\hline
\textbf{File} & \textbf{Rows} & \textbf{MCP Server} & \textbf{Primary Source} \\
\hline
\texttt{asset\_metadata.csv} & 20 & IoT & Synthetic + FDD labels \\
\texttt{sensor\_readings.csv} & 96{,}485 & IoT, TSFM & FDD time-series + Monitoring \\
\texttt{failure\_modes.csv} & 7 & FMSR & DGA labels (curated) \\
\texttt{dga\_records.csv} & 20 & FMSR & DGA + Health Index \\
\texttt{rul\_labels.csv} & 620 & TSFM & FDD RUL labels \\
\texttt{fault\_records.csv} & 506 & WO & Fault Records \\
\hline
\end{tabular}
\end{table}

\subsubsection{\texttt{asset\_metadata.csv}}

One row per transformer.  Static nameplate attributes (manufacturer, location,
voltage class, rating) are randomly sampled from representative real-world
values because none of the source datasets include nameplate data.
\texttt{health\_status} (0 = healthy, 1 = degraded, 2 = critical) is derived
from the FDD category; \texttt{rul\_days} is taken directly from the FDD \& RUL
label file.

\begin{verbatim}
transformer_id, name, manufacturer, location, voltage_class,
rating_kva, install_date, age_years, health_status,
fdd_category, rul_days, in_service
\end{verbatim}

\subsubsection{\texttt{sensor\_readings.csv}}

Long-format time-series with one row per (transformer, timestamp, sensor).
Contains 96{,}485 rows from two sub-pipelines:

\paragraph{FDD \& RUL sub-pipeline (all 20 transformers).}
The four gas columns (\texttt{H2}, \texttt{CO}, \texttt{C2H4}, \texttt{C2H2})
from each transformer's time-series file are extracted and scaled from
normalised ratios to physically plausible ppm values using empirical scale
factors derived from the DGA Fault Classification dataset
(Table~\ref{tab:scale}).  Each of the 420 time steps is assigned a synthetic
daily timestamp starting 2024-01-01.

\begin{table}[h]
\centering
\caption{Scale factors applied to FDD \& RUL normalised gas values.}
\label{tab:scale}
\begin{tabular}{lll}
\hline
\textbf{Gas} & \textbf{Scale factor} & \textbf{Resulting range (approx.)} \\
\hline
H2   & 50{,}000 & 100–200 ppm \\
CO   & 20{,}000 & 200–600 ppm \\
C2H4 & 5{,}000  & 5–25 ppm \\
C2H2 & 1{,}000  & 0.1–1 ppm \\
\hline
\end{tabular}
\end{table}

\paragraph{Real monitoring sub-pipeline (T-001 only).}
Oil temperature (OTI), winding temperature (WTI), per-phase voltage (VL1), and
per-phase current (IL1) are extracted from the Current \& Voltage Monitoring
dataset and assigned to T-001.  Zero-padded initialisation rows are dropped.
This provides T-001 with real measured temperature and electrical readings
alongside its FDD gas data, enabling richer IoT scenarios.

\begin{verbatim}
transformer_id, timestamp, sensor_id, value, unit, source
\end{verbatim}

Sensor IDs follow the convention \texttt{dga\_\{gas\}\_ppm} for dissolved-gas
readings and \texttt{\{measurement\}\_\{unit\}} for electrical readings
(e.g.\ \texttt{oil\_temp\_c}, \texttt{voltage\_l1\_v}).

\subsubsection{\texttt{failure\_modes.csv}}

A static catalogue of seven fault types, curated from the unique labels in the
DGA Fault Classification dataset and enriched with IEC 60599 standard codes,
associated key gases, severity ratings, and recommended maintenance actions.
This catalogue is used by the FMSR server to answer diagnostic queries such as
``which sensors are elevated for arc discharge?''

\begin{verbatim}
failure_mode_id, name, dga_label, description, severity,
iec_code, key_gases, recommended_action
\end{verbatim}

\subsubsection{\texttt{dga\_records.csv}}

One DGA gas-concentration snapshot per transformer, representing the most
recent analysis result.  Healthy transformers (T-001–T-010) are sourced from
the highest-scoring rows of the Health Index dataset; faulty transformers are
sourced from the DGA Fault Classification dataset, matched to the transformer's
health tier:

\begin{itemize}
  \item Minor fault (T-011–T-015): low/middle-temperature overheating and
        partial discharge records.
  \item Serious fault (T-016–T-020): arc discharge, spark discharge, and
        high-temperature overheating records.
\end{itemize}

\begin{verbatim}
transformer_id, sample_date, dissolved_h2_ppm, dissolved_ch4_ppm,
dissolved_c2h2_ppm, dissolved_c2h4_ppm, dissolved_c2h6_ppm,
dissolved_co_ppm, dissolved_co2_ppm, fault_label, source_dataset
\end{verbatim}

\subsubsection{\texttt{rul\_labels.csv}}

Daily RUL and health index estimates for each transformer over the 30-day
window (2024-01-01 to 2024-01-31).  The terminal RUL value (day 30) is taken
directly from \texttt{labels\_rul\_train.csv}; earlier days are back-filled
assuming linear degradation at one RUL-day per calendar day.  The health index
is normalised by the dataset maximum (1{,}093 days).

\begin{verbatim}
transformer_id, timestamp, rul_days, health_index, fdd_category
\end{verbatim}

\subsubsection{\texttt{fault\_records.csv}}

All 506 rows from the Smart Grid Fault Records dataset, each assigned a
\texttt{transformer\_id} via weighted random sampling.  Transformers in higher
fault tiers receive proportionally more fault events:

\begin{table}[h]
\centering
\caption{Fault record assignment weights by health tier.}
\label{tab:weights}
\begin{tabular}{ll}
\hline
\textbf{Tier} & \textbf{Relative weight} \\
\hline
Healthy (long life)  & 0.5 \\
Healthy (aging)      & 1.0 \\
Minor/moderate fault & 2.0 \\
Serious fault        & 3.5 \\
\hline
\end{tabular}
\end{table}

Column names are normalised from the original dataset's mixed-case, spaced
format to snake\_case for consistency with the other processed files.

\begin{verbatim}
transformer_id, fault_id, fault_type, location, voltage_v,
current_a, power_load_mw, temperature_c, wind_speed_kmh,
weather_condition, maintenance_status, component_health,
duration_hrs, downtime_hrs
\end{verbatim}

\subsection{Limitations and Assumptions}
\label{sec:limitations}

\begin{enumerate}

\item \textbf{Synthetic join key.}  The \texttt{transformer\_id} linkage is
constructed rather than observed.  This is a deliberate trade-off: it allows
multi-domain scenarios while being transparent about the data provenance.  We
document the full mapping in \texttt{data/build\_processed.py} so that
scenarios can be reproduced or re-stratified.

\item \textbf{Normalised gas scaling.}  The FDD \& RUL dataset stores gas
concentrations as normalised ratios rather than ppm.  The scale factors in
Table~\ref{tab:scale} were derived by aligning median FDD values to the median
DGA ppm values for the same gases.  They are physically plausible but
approximate; they do not affect the relative ordering of transformers or the
trends within a time-series.

\item \textbf{Nameplate metadata.}  No source dataset includes transformer
nameplate data (manufacturer, rating, location).  These fields are sampled from
realistic value sets and are primarily intended to support the IoT server's
\texttt{get\_asset\_metadata} tool; they have no bearing on quantitative
evaluation.

\item \textbf{Single real transformer.}  The Current \& Voltage Monitoring
dataset captures readings from one physical transformer.  Its time-series is
assigned to T-001 to enrich that unit's sensor profile.  The remaining 19
transformers do not have real electrical monitoring data.

\item \textbf{DGA healthy examples.}  The DGA Fault Classification dataset
contains only faulty examples.  Healthy DGA records are drawn from the
Transformer Health Index dataset, which uses overlapping but not identical gas
columns (no CH4 column; includes nitrogen, oxygen, DBDS, and others).  Missing
values are set to zero.

\item \textbf{RUL linearity.}  The 30-day RUL time series is generated by
linear back-fill from the label file's terminal value.  True degradation is
non-linear, but this is sufficient for the TSFM server's forecasting scenarios
during development.  When the full 420-step FDD time-series is available, the
TSFM server will use the actual sensor trajectory.

\end{enumerate}

\subsection{Reproducibility}
\label{sec:reproducibility}

The pipeline is fully deterministic given the same random seed
(\texttt{SEED=42}).  All raw data is downloaded via the Kaggle API
(\texttt{data/build\_processed.py}); raw files are gitignored.
To reproduce:

\begin{verbatim}
# 1. Set Kaggle credentials
echo '{"username":"<user>","key":"<key>"}' > ~/.kaggle/kaggle.json
chmod 600 ~/.kaggle/kaggle.json

# 2. Download raw data
python3 data/download_raw.py        # (downloads via Kaggle API)

# 3. Build processed CSVs
python3 data/build_processed.py

# 4. (Optional) generate synthetic sample data for offline testing
python3 data/generate_synthetic.py
\end{verbatim}

\noindent
A standalone synthetic generator (\texttt{data/generate\_synthetic.py})
produces structurally identical CSVs without any Kaggle download, using
normally-distributed values parameterised by health tier.  This is used for
CI testing and offline MCP server development.
